%%%%%%%%%%%%%%%%%%%%%%%%%%%%%%%%%%%%%%%%%%%%%%%%%%%%%%%%%%%%%%%%%%%%%%
% Lớp 10 | Thống kê | Các số đặc trưng đo mức độ phân tán của mẫu số liệu
% Dạng 0D6?4-3: Giá trị bất thường của mẫu số liệu
% Số câu: 1  (sinh tự động bởi tach_cau.py; từ nay thầy sửa trực tiếp file này)
%%%%%%%%%%%%%%%%%%%%%%%%%%%%%%%%%%%%%%%%%%%%%%%%%%%%%%%%%%%%%%%%%%%%%%

%%%%%%%%%%%%%%%% TRẢ LỜI NGẮN %%%%%%%%%%%%%%%%

% nguồn: MaHoaNH_TraLoiNgan_Lop10
\begin{bt}%[0D6H4-3]%[ma:da2d422e]
	Cho mẫu số liệu thống kê  sau
	\[-3 \quad 5 \quad 10 \quad 12 \quad 14 \quad 18  \quad 24 \quad 26 \quad 49 \quad 60\]
	Trong các giá trị của mẫu số liệu trên, hãy tìm giá trị ngoại lệ  của mẫu số liệu.
	 \shortans{60}
	\loigiai{ 
		Từ bảng  số liệu ta tìm được số trung vị $Q_2=\dfrac{14+18}{2}=16$, tứ phân vị thứ nhất $Q_1=10$, tứ phân vị thứ ba $Q_3=26$ và khoảng tứ phân vị $\Delta_Q=26-10=16$. Ta có \[\left[Q_1-1{,}5\Delta_Q;Q_3+1{,}5\Delta_Q\right]=\left[-14;50\right].\]
		Từ đó, ta có $60$ là số liệu ngoại lệ duy  nhất của mẫu số liệu.
	} 
\end{bt}
